\section{INTRODUÇÃO}

A nota média de um anúncio no Airbnb resume a opinião de dezenas ou centenas de hóspedes em um único número entre 1 e 5. Esse número é quase sempre alto. Em uma análise de anúncios da plataforma, \citet{zervas2021} encontraram que cerca de 95\% dos anúncios exibem média de 4,5 ou 5 estrelas. A literatura sobre avaliações online descreve distribuições concentradas no topo da escala em outros mercados \citep{hu2009,filippas2022}. Os dados deste estudo repetem o padrão: entre os 28.890 anúncios do Rio de Janeiro analisados, 9.373 (32,4\%) têm nota exatamente igual a 5,0. Uma nota tão comprimida distingue mal os anúncios entre si e diz pouco sobre o que levou o hóspede a avaliar bem ou mal.

A informação que falta à nota está no texto. Ao escrever, o hóspede menciona a limpeza, a localização, a comunicação com o anfitrião, o estado do banheiro ou o ruído da rua, em linguagem própria, frequentemente em mais de um idioma. Esse conteúdo raramente é examinado de forma sistemática por quem administra o anúncio. Para o anfitrião ou para uma empresa gestora, a pergunta útil vai além da nota: em que aspecto o anúncio se afasta dos pares e que queixas se repetem.

Este trabalho investiga, com avaliações de hóspedes do Airbnb no Rio de Janeiro (Inside Airbnb), o que o texto das avaliações revela sobre a nota. A pergunta central desdobra-se em quatro perguntas específicas. A primeira trata da medição: para quais aspectos é possível medir, de forma confiável, o sentimento do hóspede a partir de milhões de frases em português, inglês e espanhol? A segunda trata da associação com a nota: quais aspectos e quais queixas e elogios concretos mais se associam à nota global e às subnotas oficiais? A terceira trata da aplicação: como transformar esse resultado em um diagnóstico por anúncio, comparado ao bairro e ao tipo de acomodação? A quarta, de caráter secundário, trata do índice único: um índice sintético construído a partir do texto acrescenta poder explicativo à nota oficial?

\textbf{Objetivo geral.} Decompor, por aspecto e por tópico, o que o texto das avaliações de hóspedes do Airbnb no Rio de Janeiro revela sobre a nota dos anúncios, com medição validada e limites declarados e disponibilizar o resultado como um instrumento de diagnóstico por anúncio.

\textbf{Objetivos específicos.} Os objetivos seguem a ordem das perguntas. O primeiro é medir o sentimento por aspecto a partir das sentenças das avaliações e classificar os aspectos segundo a confiança da medida. O segundo é estimar a associação dos aspectos, das queixas e dos elogios com a nota global e com as subnotas oficiais. O terceiro é construir o diagnóstico por anúncio, comparado ao bairro e ao tipo de acomodação e formalizá-lo como produto técnico. O quarto é testar se um índice sintético construído a partir do texto acrescenta poder explicativo à nota oficial.

\textbf{Justificativa.} A nota concentrada no topo da escala orienta pouco quem administra um anúncio. Anfitriões e empresas gestoras precisam saber que queixas se repetem e quanto elas pesam e a plataforma não entrega essa leitura. Do ponto de vista científico, a literatura descreve a compressão das notas e os temas das avaliações, mas a validação explícita da medida por aspecto, com amostra ponderada pela população e mais de um anotador, é rara, sobretudo para avaliações em português. Do ponto de vista técnico, o estudo parte de dados públicos e produz um procedimento reprodutível, cujo resultado final é uma ferramenta de uso prático.

\textbf{Aderência à área.} O trabalho se insere no Programa de Mestrado Profissional em Gestão de Sistemas de Engenharia por tratar da gestão da qualidade de serviços e do apoio à decisão. Ele transforma dados não estruturados, as avaliações de hóspedes, em indicadores de qualidade percebida por aspecto, com método de validação e entrega um produto técnico para gestores de hospedagem. A dissertação vincula-se à área de concentração em Sistemas de Engenharia e à linha de pesquisa Modelagem Computacional do programa.

O estudo oferece quatro contribuições. A primeira é a validação da medição em camadas: frases rotuladas por um pesquisador e ponderadas pela população de sentenças, comparação de dois modelos de sentimento com regra de escolha fixada antes da execução do segundo, concordância com as subnotas oficiais, checagem com um segundo método de tópicos e rotulagem independente de 100 frases por um segundo anotador. A segunda é a decomposição da nota por aspecto e por tópico, que permite apontar queixas específicas. A terceira é o diagnóstico por anúncio, formalizado como produto técnico. A quarta é o resultado sobre o índice sintético, apresentado como resultado nulo, com suas duas leituras possíveis. Uma contribuição metodológica adicional resulta da verificação feita na base: os campos de ocupação e receita estimadas do Inside Airbnb são funções exatas do número de avaliações recentes e, por isso, só podem ser tratados como índices de volume de avaliações.

Todas as medidas são associações. O estudo não estima efeitos causais, não observa ocupação nem receita reais e se limita a um município, uma plataforma e uma coleta.

O texto está organizado em nove seções, seguidas de apêndices. A Seção 2 reúne a literatura sobre avaliações online, sobre o conteúdo das avaliações do Airbnb e sobre os métodos empregados. A Seção 3 descreve os dados e o método. A Seção 4 apresenta a validação da medição. A Seção 5 traz os resultados sobre a associação entre texto e nota, a Seção 6 a aplicação e a Seção 7 a análise do índice sintético. A Seção 8 discute os resultados e as limitações e a Seção 9 apresenta as conclusões. Os apêndices reúnem o registro datado da regra de escolha do modelo, os três relatórios de diagnóstico, as estratégias de busca bibliográfica e a tabela completa dos ganhos de $R^2$.
